\section{МЕТОДЫ ЦЕНТРАЛИЗОВАННОГО УПРАВЛЕНИЯ АВТОНОМНЫМИ АГЕНТАМИ}

% это вероятно пойдет в предыдущую главу
Важнейшей задачей управления автономными агентами является задача их 
маршрутизации --- построения глобальных маршрутов для их перемещения.


\subsection{Постановка задачи многоагентного планирования}

В общем виде задача многоагентного планирования задается графом $G = (V, E, W)$, где $V$ --- множество вершин,
$E$ --- множество ребер ($e = \left\{ v_i, v_j \right\}$),  $W : E \rightarrow \mathbb{R}$ --- функция,
задающая вес ребра и множеством агентов $A = \left\{ 1, \ldots, n \right\}$.
Граф может быть как ориентированным, так и нет, в данной работе исходные данные представляют из себя
неориентированные планарные графы $W : E \rightarrow ||E||_2$ % (CHECK)

В каждый момент времени для агента $i$ определены начальная (стартовая) $s_i$ и конечная (целевая) $g_i$ вершины графа.
Эти вершины могут совпадать, что позволяет агентам ожидать следующую цель.
Множество агентов считается неизменным. Агенты перемещаются по графу с постоянной скоростью $v$.

Задача многоагентного планирования состоит в нахождении путей в графе между начальными и целевыми вершинами каждого из агентов,
которые будут удовлетворять некоторым ограничениям, которые будут определены позднее.

Метрики

% \subsection{Multi-agent path finding}

% \subsubsection{Conflict based search}

